% Nouveau EGA V (cote 25): archive cote 25, site page 100 (PDF page 101).
% typescript. Draft compared: none.
% AI-transcribed by Claude Opus 5.5 (claude-opus-5-5, adaptive thinking, effort high), Message Batches API, batch ega25b (msgbatch_01A7n5YuXubUPNK1zsepi8RG), 2026-11-13; not yet checked by eye.
\scanpage{25/100}
\mspage{3}

% NOTE: a single long handwritten diagonal stroke crosses the whole page from top to bottom; it is transcribed as a deletion of each paragraph
\deleted{Prouvons l'existence des $(S_n)$, $\underline{M}$ . Grâce à l'unicité déjà prouvée,
la question est \uline{locale} sur S, i.e. on est ramené à prouver le
\uline{Corollaire} 1.5. Soit $\underline{E}$ loc libre de rang fini sur S, P=P($\underline{E}$), $\underline{L}$ un
Module inversible sur P, s $\in$ S, alors il existe un voisinage ouvert U de s
et un entier n $\in \underline{Z}$ , tel que $\underline{L}\,|\,f^{-1}(U)$ soit isomorphe à $\underline{O}_P(n)\,|\,f^{-1}(U)$.}

\deleted{\hand{\unsure{[}}Bien entendu, lorsque le rang de E en s est $\geq$ 2, l'entier n est bien
déterminé. D'ailleurs 1.5. est trivial lorsque le rang de $\underline{E}$ en s est $\leq$ 1\hand{\unsure{]}}.
Noter que $\underline{E}$ est de la frime, car la question étant locale, on peut supposer
déjà $\underline{E} = \underline{O}_S^{r+1}$, donc P = $P_S^{r\xout{+1}}$\hand{\unsure{]}}. Par le procédé bréveté du par.8, on est
ramené de plus au cas où S est noethérien. Nous précédons endeux pas:}

\deleted{a) S est le spectre d'un corps k. Nous savons que $\underline{L}$ est défini
par un Module gradué de type fini L sur l'anneau gradué k$[t_0,\ldots,t_r]$.
On sait aussi que la restriction \inserted{$L'$} de $\tilde{L}$ à l'espace affine épointé
$E_k^{r+1}$ - (0) n'est autre que l'image inverse de $\underline{L}$ par le morphisme de
projection canonique \inserted{q:} $E_k^{r+1}$ - (0) $\longrightarrow P_k^r$ , et c'est donc un Module inver-%
sible. Soit $i : E^{r+1}$-(0)$\to E^{r+1}$ l'immersion canonique, il résulte
alors du fait que l'anneau \inserted{affine} k$[t_0,..,t_r]$ de $E^{r+1}$ est factoriel, donc à
fortiori son localisé en le point 0 de $E^{r+1}$ est factori\xout{\unsure{e}}el, et du fait
que ce dernier anneau est de dimension $\geq$ 2, que $i_*(\tilde{L'})$ est un Module
inversible, \xout{\illegible{1}} donc correspond à un module inversible M sur
k$[t_0,..,t_r]$. D'ailleurs M = $\Gamma(E^{r+1}$-0, $\unsure{q'|q^*}(\underline{L}))$ est gradué de façon natu-%
relle, enfin l'homomorphisme L $\to$ M est évidemment un isomorphisme
entous les \xout{\illegible{2}} points de $E^{r+1}$ distincts de 0. Donc quitte à
remplacer L par M, on est ramené au cas où L est inversible. Mais
k$[t_0,..,t_r]$ étant factoriel, L est donc \uline{libre} de rang 1 , d'abord
enfaisant abstraction de sa graduation; mais un lemme standard (qui se
trouve sans doute dans Bourbaki) implique qu'il est alors libre de rang 1
entant que module gradué, ce qui implique que le Module associé sur P}

\mspage{100}
